% GENERATED FILE. Edit canonical lessons, then run npm run generate:books.
% canonical-source-hash: fnv1a64:89ffa688255fddbc
% canonical-lessons: TA-C14-kaalangal

\chapter{Seasons}
\label{ch:ta-seasons}

This chapter is generated from the canonical micro-lessons used by Language Ladder.
Each section stays independently resumable and preserves the authored prerequisite order.

\section[\ta{வசந்த} \ta{காலம்} \ta{கோடைக்} \ta{காலம்} \ta{மழைக்} \ta{காலம்} \ta{குளிர்} \ta{காலம்}]{\ta{வசந்த} \ta{காலம்}, \ta{கோடை}, \ta{மழை}, \ta{குளிர்} — the four words, and the real season underneath}

\label{lesson:TA-C14-kaalangal}



\begin{quote}
\textbf{Warm-up.} {[}PAUSE 2s{]} Tamil has words that line up with Spanish's four seasons — but, just like the Hindi lesson on \dv{ऋतु}, forcing them into a clean four-way split hides what actually matters on the ground.
\end{quote}

\subsection*{You'll want to know: The four words}
\noindent
\begin{tabularx}{\linewidth}{@{}>{\raggedright\arraybackslash}X>{\raggedright\arraybackslash}X>{\raggedright\arraybackslash}X@{}}
\toprule
\textbf{Tamil} & \textbf{meaning} & \textbf{source} \\
\midrule
\textbf{\ta{வசந்த} \ta{காலம்}} (\emph{vasantha kaalam}) & "spring" & \emph{vasantham} is a \textbf{Sanskrit} loan \\
\textbf{\ta{கோடைக்} \ta{காலம்}} (\emph{kodai kaalam}) & "summer" & \textbf{\ta{கோடை}} \emph{kodai} is \textbf{native} Tamil for "heat" \\
\textbf{\ta{மழைக்} \ta{காலம்}} (\emph{mazhai kaalam}) & "rainy season" & \textbf{\ta{மழை}} \emph{mazhai}, native Tamil for "rain" \\
\textbf{\ta{குளிர்} \ta{காலம்}} (\emph{kulir kaalam}) & "winter/cold season" & \textbf{\ta{குளிர்}} \emph{kulir}, native Tamil for "cold" \\
\bottomrule
\end{tabularx}

Every one of these is built the same way: a word for the \textbf{feel} of the season (heat, rain, cold — or the borrowed \emph{vasantham}) plus \textbf{\ta{காலம்}} (\emph{kaalam}), "\textbf{time/season}." Be honest about \emph{kaalam} itself, too: it's ultimately from Sanskrit \textbf{\dv{काल}} (\emph{kāla}, "time") — but it was assimilated into Tamil so long ago and so thoroughly that it functions as completely ordinary vocabulary today, quite unlike a fresh, still-foreign-feeling loan like \emph{vasantham}.

\begin{culture}
Just as Hindi's traditional calendar makes the monsoon (\emph{varṣā}) its own central season rather than a footnote, Tamil Nadu's actual climate is shaped far more by \textbf{\ta{மழைக்காலம்}} (\emph{mazhai kaalam}, "the rains") than by any crisp four-way Western split. There's no dramatic autumn leaf-fall in a tropical climate, and "spring" as a distinct, keenly-felt season is a much weaker idea here than in temperate Europe. The four words above are real and used — but the rains, not a European four-season year, are what actually organizes agricultural and everyday life across much of Tamil Nadu.

(Note for the Tamil-speaking reviewer: please check whether these are the words and register you'd teach first, and whether there's more to say about how Tamil Nadu's own seasonal year is actually divided in practice.)
\end{culture}

\subsection*{Guided Practice}
{[}PAUSE 1s{]}

\begin{itemize}
\raggedright
  \item {[}YOU SAY: "kodai" — heat/summer, native — then "kaalam," time/season{]}
  \item {[}YOU SAY: "mazhai kaalam" — the rains, the real central season{]}
  \item {[}YOU SAY: "kulir kaalam" — the cold season{]}
  \item {[}YOU SAY: the one loanword — "vasantha kaalam," spring, from Sanskrit{]}
\end{itemize}

\begin{tcolorbox}[breakable,colback=teal!4,colframe=teal!35!black,title={Wrap-up recall}]
{[}PAUSE 3s{]} Which of the four season-words is a Sanskrit loan, and which are native Tamil? (\textbf{\ta{வசந்தம்}} \emph{vasantham} is Sanskrit; \textbf{\ta{கோடை}, \ta{மழை}, \ta{குளிர்}} are all native.) What season actually shapes life in Tamil Nadu the most, and why doesn't the four-season frame capture that? (\textbf{\ta{மழைக்காலம்}}, the \textbf{rains} — Tamil Nadu's tropical climate doesn't have a sharp four-way split the way temperate Europe does.)
\end{tcolorbox}
